\section{Evaluation}
\label{sec:evaluation}

This section separates two kinds of evidence. First, we report measurements that are reproducible directly from the current database artifacts: logical database sizes, preprocessed table sizes, Merkle proof material, OnionPIR stores, and exact two-server DPF communication per PBC round. Second, we define the controlled benchmark matrix needed for latency and throughput claims. The distinction matters because artifact size and DPF wire cost are deterministic functions of the database layout, while query latency depends on hardware, cache state, network placement, and server concurrency.

\subsection{Current Artifact Measurements}

The current measurement snapshot uses a full Bitcoin UTXO checkpoint at height $948{,}454$ and the delta database from height $940{,}611$ to $948{,}454$. Measurements were generated by \texttt{benchmarks/collect\_artifact\_metrics.py} from the local artifact root \texttt{/Volumes/Bitcoin/data}. The script parses the cuckoo-table headers, records artifact sizes, computes DPF domain sizes, and emits both CSV files and the tables included below. The measured machine was an Apple M2 Ultra system with 24 logical cores and 64~GiB of memory; exact machine and git provenance are recorded in \texttt{benchmarks/results/environment.json}.

% Auto-generated by benchmarks/collect_artifact_metrics.py.
\begin{table}[ht]
\centering
\small
\begin{tabular}{lrr}
\hline
Artifact & Full snapshot & 7,843-block delta \\
\hline
Raw UTXO/delta input & 10.44 GiB & 453.8 MiB \\
Logical index & 1.25 GiB & 120.0 MiB \\
Logical chunk payload & 3.02 GiB & 324.5 MiB \\
DPF/Harmony index cuckoo & 2.06 GiB & 198.2 MiB \\
DPF/Harmony chunk cuckoo & 10.49 GiB & 1.10 GiB \\
Bucket Merkle artifacts & 4.36 GiB & 452.9 MiB \\
OnionPIR NTT store & 14.48 GiB & 1.77 GiB \\
OnionPIR index store & 12.30 GiB & 1.10 GiB \\
\hline
\end{tabular}
\caption{Measured artifact sizes for the current Bitcoin PIR full snapshot at height 948,454 and the delta from 940,611 to 948,454.}
\label{tab:artifact-sizes}
\end{table}

\begin{table}[ht]
\centering
\small
\begin{tabular}{llrrrr}
\hline
Database & Level & Bins/table & $n$ & Key & Round traffic \\
\hline
Full & INDEX & 567,558 & 20 & 355 B & 119.2 KiB \\
Full & CHUNK & 1,066,928 & 21 & 371 B & 157.2 KiB \\
Delta & INDEX & 53,287 & 16 & 290 B & 100.2 KiB \\
Delta & CHUNK & 112,344 & 17 & 307 B & 137.2 KiB \\
\hline
\end{tabular}
\caption{Exact two-server DPF communication per PBC round, derived from current cuckoo-table headers. Each round sends two DPF keys per group to each server and receives one bucket response for each key.}
\label{tab:dpf-round-costs}
\end{table}



Several immediate conclusions follow from Table~\ref{tab:artifact-sizes}. The DPF/Harmony chunk cuckoo table dominates the two-server storage footprint because each chunk bucket stores inlined chunk payloads. OnionPIR shifts this cost into the NTT-preprocessed store and the indexed FHE layout. Delta databases are much smaller than full checkpoints: for the $940{,}611 \to 948{,}454$ interval, the logical chunk payload is $324.5$~MiB compared with $3.02$~GiB for the full checkpoint, and the DPF/Harmony chunk cuckoo table is $1.10$~GiB compared with $10.49$~GiB.

Table~\ref{tab:dpf-round-costs} gives the exact two-server DPF communication for one PBC round at the INDEX and CHUNK levels. These are not extrapolated timing measurements: they are derived from the current table headers. The full checkpoint uses $n=20$ for INDEX and $n=21$ for CHUNK, while the measured delta uses $n=16$ and $n=17$. The smaller delta domain lowers DPF key size, though the bucket response size remains fixed by the slot layout.

\subsection{Latency and Throughput Benchmark Matrix}

The remaining evaluation should benchmark the same implementation under multiple scales and environments. A publishable latency table should not use a single full-scale run alone: it should show the trend as the database grows and separate preprocessing costs from online query costs.

\begin{table}[ht]
\centering
\small
\begin{tabular}{p{0.22\linewidth}p{0.70\linewidth}}
\hline
Axis & Required coverage \\
\hline
Database scale & Synthetic smoke databases for correctness and overhead isolation; the existing $940{,}611 \to 944{,}000$ delta; the existing $940{,}611 \to 948{,}454$ delta; the full $948{,}454$ checkpoint; and at least one future checkpoint to show growth with the Bitcoin UTXO set. \\
Protocol path & DPF-PIR INDEX and CHUNK rounds; HarmonyPIR hint generation, hint download, and online query; OnionPIR key generation, query generation, server answer, and decryption; Merkle proof retrieval and verification for all protocol paths. \\
Environment & Local single-machine warm-cache runs; local cold-cache runs after page-cache eviction or reboot; two-host LAN/WAN WebSocket runs; browser/WASM client runs; and SEV-SNP deployment runs when the TEE-hardened path is evaluated. \\
Metric & Artifact size, preprocessing time, server resident set size, p50/p95/p99 online latency, sustained queries per second, client upload/download bytes, client CPU time, and verification failure rate over randomized test queries. \\
Repetition & At least 30 online-query trials per configuration after a fixed warmup, plus independent rebuild or reload trials for preprocessing measurements. Report median and tail latency, not only mean. \\
\hline
\end{tabular}
\caption{Benchmark matrix required before making final latency and throughput claims.}
\label{tab:benchmark-matrix}
\end{table}

For DPF-PIR, the primary scaling variable is the number of bins per group. The benchmark should report INDEX and CHUNK separately because the INDEX response is small but drives chunk discovery, while CHUNK scans dominate payload retrieval. The current exact communication table gives the wire cost; the timing benchmark must additionally measure DPF key generation, per-server scan time, and client-side XOR/tag verification.

For HarmonyPIR, offline and online measurements must be separated. The offline hint phase is paid when the client prepares a stateful session and should be reported as a one-time cost per database and per PRP backend. The online phase should report request construction, server response time, response processing, and query budget consumption. The paper should include both HMR12 and FastPRP results if both are supported by the deployment being evaluated.

For OnionPIR, the evaluation must distinguish database preprocessing from online queries. The NTT store is large but amortized across all clients. Per-session client key generation and key upload should be reported separately from per-query upload, server answer time, response size, and decryption time. Because the OnionPIR benchmark currently includes a self-contained random-data path, final paper numbers should either run against the current real artifacts or explicitly label the results as synthetic-capacity measurements.

\subsection{Reproducibility Artifacts}

All benchmark-backed claims in this section should be regenerated from committed artifacts rather than copied from terminal output. The current repository records:
\begin{itemize}[nolistsep,leftmargin=*]
    \item \texttt{benchmarks/results/artifact\_metrics.csv}: logical and preprocessed artifact sizes for the measured databases.
    \item \texttt{benchmarks/results/dpf\_round\_costs.csv}: exact DPF domain and communication costs derived from cuckoo headers.
    \item \texttt{benchmarks/results/environment.json}: hardware, operating system, and git revision metadata for the collection run.
    \item \texttt{benchmarks/results/artifact\_tables.tex}: generated tables included in this paper.
\end{itemize}

The next benchmark step is to add latency runners that emit machine-readable CSV with one row per trial. Those runners should be responsible for warmup, trial count, cache-state labeling, and protocol-specific phases, so the paper can report medians and tail latencies without hand-editing results.
